\clearpage
\fancyhead[L]{\small BODE. ASTRONOMISCHES JAHRBUCH.}
\hypertarget{native-553-lower}{}\label{native-553-lower}
\begin{center}\small Göttingische gelehrte Anzeigen. Stück 22. Seite 209..220. 1813 Februar 6.\\[3pt]\rule{.72\linewidth}{.3pt}\end{center}

\emph{Astronomisches Jahrbuch für das Jahr} 1814, nebst einer Sammlung der neuesten in die astronomischen Wissenschaften einschlagenden Abhandlungen, Beobachtungen und Nachrichten, berechnet und herausgegeben von J. E. \textsc{Bode}, königl. Astronom und Miglied der Akademie der Wissenschaften. 1811. Bei dem Verfasser, und in Commission bei J. E. \textsc{Hitzig}. Berlin. 275 Seiten in Octav, u. Kupfertafel.

Astronomisches Jahrbuch für das Jahr 1815 u. s. w. Eben daselbst 1812. Gleichfalls 275 Seiten in Octav und eine Kupfertafel.

Ohne die Bemerkungen zu wiederholen, welche wir über die Einrichtung des astronomischen Kalenders bereits bei der Anzeige früherer Bände dieses Jahrbuchs gemacht haben, schränken wir uns hier blos auf die Anzeige der angehängten Abhandlungen ein. Der Jahrgang von 1814 enthält folgende 34 Artikel. 1) Die mittlere astronomische Strahlenbrechung, nach \textsc{Laplace}. 2) Astronomische Beobachtungen in den Jahren 1809 und 1810 auf der königl. Sternwarte zu Kopenhagen angestellt von \textsc{Bugge}, betreffen die Ceres, Vesta, den Uranus, Saturn, Jupiter, Verfinsterungen von dessen Trabanten, und Sternbedeckungen. 3) Ueber eine Methode, die Zeit zu bestimmen durch Messung einer Distanz der Sonne von einem festen und bekannten Punkt am Horizonte, von Hrn. J. F. \textsc{von Beek Calkoen}. Dies \SourceLeaf{554}Methode kann, bei gehöriger Vorsicht, brauchbare Resultate geben: allein sie muss nicht auf Gegenstände am Horizonte beschränkt, und die Refraction und Parallaxe nicht vernachlässigt werden. Man kann auf beides gehörig Rücksicht nehmen, ohne das Verfahren sehr weitläufig zu machen. 4) Astronomische Beobachtungen, auf der kaiserl. Sternwarte zu Wien angestellt von Hrn. Doctor \textsc{Triesnecker} und Hrn. Professor \textsc{Bürg}, enthalten Trabantenverfinsterungen, Sternbedeckungen, Planeten- und Sonnenbeobachtungen. 5) Ideen zur Perturbationsrechnung, nach \textsc{Kepler}, nebst Anmerkungen, von Hrn. J. W. \textsc{Pfaff}, verdienen zum Theil beachtet zu werden, obwohl uns Manches unreif oder übereilt scheint. 6) Astronomische Beobachtungen, auf der königl. Sternwarte zu Prag angestellt im Jahre 1810 von Hrn. Prof. \textsc{David} und Hrn. Adjunct \textsc{Bittner}. Der Art, wie hier und in den vorhergehenden Bänden dieses Jahrbuchs beobachtete Sternhöhen zu einer vermeintlichen Correction der Refractionstafeln angewandt sind, können wir unsern Beifall nicht geben. Wer über diesen delicaten Punkt arbeiten will, muss zuerst die Polhöhe seines Beobachtungsortes fester begründen, und die Beobachtungen überhaupt weit mehr vervielfältigen, als es Hr. \textsc{David} gethan hat, und darf sich nicht begnügen, die Stern-Declinationen von andern Astronomen zu entlehnen. 7) Aus einem Schreiben des Hrn. Prof. \textsc{Littrow} aus Kasan. Es ist erfreulich, daraus das Aufblühen der Astronomie an der Grenze von Asien zu ersehen. 8) Genauere Bestimmung der Lichtänderungs-Periode des Sterns $\eta$ Antinoi, vom Hrn. Prof. \textsc{Wurm}. Aus eignen Beobachtungen, welche einen Zeitraum von 24 Jahren umfassen, findet Hr. \textsc{Wurm} die Periode $7{,}1761$ Tage. Man hätte hier eine weniger willkührliche Art, die Beobachtungen zu combiniren, etwa nach der Methode der kleinsten Quadrate, und nachher eine Vergleichung mit den einzelnen Beobachtungen, wünschen können. 9) Beobachtungen des Cometen von 1807, und der totalen Sonnenfinsterniss vom 16. Juni 1806, zu Salem in den vereinigten Nordamerikanischen Staaten, von Hrn. \textsc{Bowditch}. Aus dem letzten Bande der \emph{Memoirs of the American Academy of Arts and Sciences}. Interessant ist besonders die Erzählung der die totale Sonnenfinsterniss begleitenden Umstände. 10) Nachricht von der Mannheimer Sternwarte, vom Hrn. Staatsrath \textsc{Klüber}, Curator der Sternwarte. Das hier Gesagte ist durch die seitdem erschienene besondere Schrift des Hrn. \textsc{Klüber} bereits umständlicher bekannt. 11) Astronomische Beobachtungen, auf der königl. Sternwarte zu Berlin angestellt im Jahre 1810. Wir heben aus ihnen die Bestimmung der Polhöhe der Berliner Sternwarte mit dem zweifussigen Troughtonschen Kreise aus, welche Hr. \textsc{Bode} im Mittel $52^\circ\,31'\,15''$ findet. 12) Ueber den Cometen von 1795, vom Hrn. Dr. \textsc{Olbers}. Die Beobachtungen dieses Cometen sind sehr dürftig, und die von verschiedenen Astronomen herausgebrachten Elemente weichen beträchtlich von einander ab. Hr. Dr. \textsc{Olbers} gründete eine neue Bestimmung der Elemente theils auf eine früher noch nicht benutzte Beobachtung \textsc{Herschel}'s, theils auf eine von neuem reducirte Beobachtung von Hrn. \textsc{Bode}. 13) Resultate einer Untersuchung über die Lage der Ebene des Saturnsringes, die Theorie des vierten Satelliten, die Massen des Planeten und des Ringes, und beobachtete Sternbedeckungen, von Hrn. Prof. \textsc{Bessel}. Die ausführliche Untersuchung über den vierten Saturnstrabanten ist bereits im dritten Hefte des Königsberger Archivs erschienen; die wichtigen Resultate verdienen die Aufmerksamkeit aller Astronomen. 14) Berechnung der Bahn des Cometen von 1810, von demselben. 15) Astronomische Beobachtungen, zu Kremsmünster im Jahre 1809 und 1810 angestellt von Hrn. \textsc{Derfflinger}. Bei dem Eintritte des Aldebaran am 18. September 1810 schien Hrn. \textsc{Derfflinger} der Stern ein paar Secunden hindurch mit verändertem Lichte auf dem hellen Mondsrande zu verweilen. Eben dies bemerkte auch \textsc{Bugge}, \textsc{Bode} und \textsc{David}, und einige Astronomen haben dies als etwas Ausserordentliches betrachtet. Auf der hiesigen Sternwarte, wo dieselbe Bedeckung gleichfalls beobachtet wurde, ist nichts der Art bemerkt. Wir lassen es daher dahin gestellt sein, ob dies etwas Anderes, als Irradiation des Mondsrandes bei vielleicht schon etwas ermüdetem Auge gewesen sei. 16) Beobachtung und Berechnung der Bedeckung des Aldebaran vom Monde am 18. September 1810, \SourceLeaf{555}zu Dorpat angestellt vom Hrn. Prof. \textsc{Knorre}. Von dem eben erwähnten Phänomen ist hier nichts gesagt. 17) Ueber das Höhenmessen vermittelst des Barometers, vom Hrn. Prof. \textsc{Benzenberg}. Dieser an sich lobenswerthe Versuch, diesen Gegenstand populär darzustellen, ist hier wohl nicht ganz an seinem Platze. 18) Beobachtungen der Juno und Vesta im Jahre 1811 auf der kaiserl. Sternwarte zu Wilna, von Hrn. Prof. \textsc{Sniadecki}. 19) Aus einem Schreiben des Hrn. Dr. \textsc{Pansner} in Petersburg. Dieser Artikel enthält Nachrichten von einer trigonometrischen Vermessung der Küste des Finnischen Meerbusens, deren Vollendung sehr zu wünschen ist. 20) Ueber die Genauigkeit des \textsc{Baumann}'schen Verticalkreises, von Hrn. Dr. \textsc{Pottgiesser} in Eberfeld: ein mit vieler practischer Einsicht geschriebener Aufsatz. Ausser einer Schätzung der Fehler aus der Beschaffenheit des Instruments selbst, hat der Verfasser die Genauigkeit desselben auch durch einige wirkliche Beobachtungen geprüft, welche sehr vortheilhaft dafür sprechen. Wir hätten nur gewünscht, dass dieselben zahlreicher und mannigfaltiger wären, da neuere Erfahrungen einige Bedenklichkeiten gegen die Kreise mit stehenden Säulen angeregt haben. 21) Längen- und Breitenbestimmungen einiger Oerter im Oestreichischen, nebst beobachteten Sternbedeckungen, von der Frau Reichsfreiin von \textsc{Matt}. 22) Beobachtungen über die jährliche Parallaxe von $\alpha$ Leyer, von Hrn. \textsc{Calandrelli} in Rom. Es wird hier eine jährliche Parallaxe von $5''8$ gefolgert, gegen deren Zuverlässigkeit sich aber doch noch Mehreres erinnern liesse. Der Sector, womit die Beobachtungen angestellt sind, wurde nicht umgewandt, sondern vorausgesetzt, dass sein Collimationsfehler das ganze Jahr hindurch unveränderlich geblieben sei. 23) Entwurf einer Sonnenuhr, welche die zwölfte Mittagsstunde mittlerer Zeit angibt. Durch 36 Punkte wird die in Form einer 8 geschlungene Curve gezeichnet, welche das Ende des Schattens des Zeigers auf einer Verticaluhr im mittlern Mittage in den verschiedenen Monaten des Jahrs bildet. Der Gebrauch der mittlern Zeit bei öffentlichen Uhren hat allerdings viel für sich; da indessen die Personen, welchen die Stellung der Uhren auf dem Lande oder an kleinern Oertern obliegt, selten darüber gehörig unterrichtet sind, so wird dabei der grosse Vortheil der \emph{Uebereinstimmung} der öffentlichen Uhren, an welcher doch besonders auf Poststrassen viel gelegen sein, und die beim Gebrauch der wahren Zeit leichter erhalten werden kann, leicht verloren. 24) Astronomische Ortsbestimmungen, vom Hrn. Oberprediger \textsc{Fritsch} in Quedlinburg. Hr. \textsc{Fritsch} bestimmte auf einer Reise nach Schlesien im Sommer 1810 die Polhöhe von Ballenstedt, Werningerode, Oschatz, Bischofswerda und mehrerer Punkte in Schlesien. 25) Ueber den Cometen vom Jahre 1811 und dessen Wiederkunft im August, von Hrn. Dr. \textsc{Olbers}. 26) Astronomische Beobachtungen und Bemerkungen, vom Hrn. Prof. \textsc{Gauss}, enthalten Beobachtungen der Pallas und Vesta, und Berichtigungen einiger Bemerkungen des Hrn. Dr. \textsc{Treisnecker} im Jahrbuche für 1813. Die drei folgenden Artikel geben die Ephemeriden für die Pallas, Juno und Vesta von den Herren \textsc{Nicolai}, \textsc{Wachter} und \textsc{Gerling}. 30) Beobachtungen und Elemente der Bahn des Cometen von 1811, und Zusatz zu der Theoria motus corporum coelestium vom Hrn. Prof. \textsc{Gauss}. In dem \emph{Zusatz} bemerken wir einen Druckfehler, indem S. 257 Z. 5 statt $z$ gelesen werden muss $-z$. 31) Beobachtungen des Cometen von 1811, die Elemente seiner Bahn, und Sternbedeckungen, von Hrn. Prof. \textsc{Bessel}. 32) Aus einem Schreiben des Hrn. Dr. \textsc{Koch} in Danzig. Bemerkungen über einige Fixsterne, und Bestimmung der Abweichung der Magnetnadel in Danzig ($13^\circ\,48'$). 33) Beobachtungen des Cometen von 1811 auf der königl. Sternwarte in Berlin. Bei der beigefügten Zeichnung der wahren Bahn des Cometen muss bemerkt werden, dass sie keine Projection ist. --- Unter den kleinen astronomischen Notizen, welche den letzten Artikel ausmachen, heben wir nur die Nachricht aus einem Schreiben des Hrn. Dr. B\dots\ (\textsc{Brinkley}?) aus, welcher mit einem achtfussigen Kreise die jährliche Parallaxe von $\alpha$ Leyer $2''52$ gefunden hat, welches noch nicht die Hälfte von dem oben angeführten Resultate \textsc{Calandrelli}'s ist.

Der Jahrgang 1815 liefert 31 Artikel. Den Anfang machen Bemerkungen über des Hrn. Prof. \SourceLeaf{556}\textsc{Gauss} Theoria motus corporum coelestium, von Hrn. Prof. \textsc{Littrow}. Sie beziehen sich auf die Aufgabe, aus drei geocentrischen Oertern eines Planeten seine Bahn zu bestimmen, wo Hr. \textsc{Littrow} für den freilich sehr seltenen, Fall, dass drei Hypothesen nicht ausreichen, im Gange der Rechnung der vierten Hypothese eine kleine Abänderung macht. Diese besteht darin, dass er durch denselben Kunstgriff, wodurch \textsc{Gauss} die neuen Werthe von $P$ und $Q$ bestimmt, sogleich $r, r', f, f', f''$ berechnet. Obgleich dies Verfahren zuweilen einigen Vortheil geben kann, so möchten wir es doch nicht unbedingt anrathen. Man ist im Allgemeinen nicht berechtigt, davon genauere Resultate zu erwarten, als von dem in der Theoria angewandten, und die Verkürzung des Weges wird dadurch zum Theil wieder aufgehoben, dass man die Lage der Ebene der Bahn dadurch nicht mit erhält, auf welche eben jenen Kunstgriff auch sofort anzuwenden doch zuweilen etwas misslich sein kann. Die hierauf folgende, von demselben Astronomen angegebene, Methode zur Bestimmung einer Kreisbahn aus zwei geocentrischen Oertern, halten wir für sehr zweckmässig, nur glauben wir nicht, dass es vortheilhaft sei, sie zur vorläufigen Bestimmung der Abstände bei \emph{Cometenbahnen} anzuwenden. Noch finden wir von dem Verfasser eine leichte Deduction der Aberration für Länge und Breite, und ein paar Reihenausdrücke für die Refraction, wovon der eine nicht genug erklärt ist, um darüber ein Urtheil fällen zu können. 2) Astronomische Beobachtungen, zu Pisa, Mailand und Padua angestellt (aus dem XV. Bande der Memorie di Matematica e di Fisica della Società Italiana ausgehoben). 3) Astronomische Nachrichten und Bemerkungen, physische Beobachtung des grossen Cometen von 2811, geographische Bestimmungen u.\,s.\,w. von Hrn. Prof. \textsc{Huth} in Dorpat. 4) Beobachtung des grossen Cometen von 1811, Berechnung seiner Bahn, Sternbedeckungen u.\,s.\,w. von Hrn. Prof. \textsc{Bessel}. 5) Ueber die Entdeckung eines neuen Cometen im November 1811, Beobachtungen desselben und des grossen Cometen von 1811, Beobachtung der Pallas von Hrn. Dr. \textsc{Olbers}. 6) Beobachtungen des Cometen von 1807 und die Elemente seiner Bahn, von Hrn. \textsc{Niccolo Cacciatore}, Gehülfen bei der königl. Sternwarte zu Palermo: ein sehr schätzbarer, bisher auf dem festen Lande noch. nicht bekannter, Nachtrag für die scheinbare Bahn jenes Cometen. 7) Beobachtete und berechnete Gegenscheine des Mars, der Vesta, des Jupiter und Saturn, Jupiterstrabanten-Verfinsterungen, Sternbedeckungen, im Jahre 1811, Elemente des Cometen von 1810, Beobachtung des grossen Cometen von 1811, und Berechnung seiner Bahn, von Hrn Dr. \textsc{Triesnecker}. 8) Beobachtungen des Uranus, Saturn, Mars, der Ceres, Jupiterstrabanten-Verfinsterungen, Sternbedeckungen und des grossen Cometen von 1811, auf der kaiserl. Sternwarte zu Wilna angestellt von Hrn. Prof. \textsc{Sniadecki}. 9) Astronomische Nachrichten und Beobachtungen, geographische Ortsbestimmungen u.\,s.\,w. von Hrn. Prof. \textsc{Jabbo Oltmanns}. Beobachtete Sternbedeckungen in Amsterdam und Greenwich, Beobachtungen in Südamerika, Ortsbestimmungen in Spanien und Portugal. 10) Andenken an den \textsc{Halley}ischen Cometen, von Hrn. Prof. J. W. \textsc{Pfaff} in Nürnberg. Der Verfasser überlässt sich hier der Aussicht auf die vielfachen Aufschlüsse, welche die wiederholt beobachtete Wiederkehr der Cometen einst geben wird. Freilich werden erst noch Jahrhunderte vergehen müssen, ehe auch nur ein Theil dieser Hoffnungen sich wird realisiren können. 11) Astronomische Beobachtungen, auf der königl. Sternwarte zu Berlin angestellt im Jahre 1811. 12) Beobachtungen über Jupiterstrabanten-Verfinsterungen, Mondfinsternisse, Sternbedeckungen, Refraction, Planeten-Gegenscheine und den grossen Cometen von 1811, auf der Prager Sternwarte angestellt von Hrn. Prof. \textsc{David}. 13) Beobachtungen der Pallas, des zweiten Cometen von 1811 und Elemente seiner Bahn, Sternbedeckungen u.\,s.\,w. von Hrn. Prof. \textsc{Gauss}. 14) Ueber das Zusammentreffen der Erde und des Mondes an einem und demselben Orte. Vorausberechnung der Fälle, wo die Erde einige Stunden nachher an den Platz kommt, den vorher der Mond eingenommen hatte, und umgekehrt, für die Jahre 1812--1815, Lichtenberg hatte bekanntlich einmal die Frage aufgeworfen, ob sich vielleicht im ersteren Falle auffallende Witterungsveränderungen zeigten. Wir gestehen, dass \SourceLeaf{557}wir (eben so wie Hr. \textsc{Bode} selbst) nichts davon erwarten, und bemerken nur noch, dass die Rechnung selbst sich auf die unzulässige Voraussetzung der absoluten Ruhe des Sonnensystems gründet. 15) Beobachtungen der Ceres, und der Gegenscheine des Uranus und Mars im Jahre 1811, auf der Sternwarte zu Kremsmünster angestellt von Hrn. \textsc{Derfflinger}. 16) Astronomische Bemerkungen, von Hrn. Dr. \textsc{Lamberti} in Dorpat. 17) Ueber die scheinbare Bahn des grossen Cometen von 1811, nebst einer Zeichnung. 18) Ueber die Bewegung des Doppelsterns 61 im Schwan, von Hrn. Prof. \textsc{Bessel}. Ueber diesen merkwürdigen Gegenstand ist bereits an mehreren andern Orten, und zuerst in unsern Blättern (1812 April 25 St. 67 [S. 349 d. B.]) Nachricht gegeben. 19) Verbesserung der Bestimmung der Polhöhe von Riga, von Hrn. Prof. \textsc{Sandt}. 20) Astronomische Beobachtungen, zu Paris und Greenwich angestellt in den Jahren 1805--1809, mitgetheilt an Hrn. Prof. \textsc{Oltmanns}. 21) Nachricht von sehr vollkommenen Parallelspiegeln, die vom Hrn. Mechanicus \textsc{Duve} in Berlin verfertigt werden, von Hrn. Prof. \textsc{Fischer} in Berlin. Die Wichtigkeit des Umstandes, dass bei Reflexions-Instrumenten die Glasspiegel vollkommen ebene und parallele Flächen haben, macht die Bemühungen des Hrn. \textsc{Duve}, welche nach Hrn. \textsc{Fischer} wohl gelungen sind, sehr schätzbar. Bemerkenswerth ist die hier erzählte Erfahrung, dass ein solcher Parallelspiegel seine vorige Vollkommenheit verloren hatte, als durch einen Zufall ein kleines Stück davon abgebrochen war. 22) Astronomische Beobachtungen und Bemerkungen, von Hrn. Joseph \textsc{Bayr} zu Kloster Hradisch bei Ollmütz. Geographische Bestimmung von Hradisch, und Polhöhe von Troppau. 23) Ueber den Einfluss der \textsc{Dalton}'schen Theorie auf das Höhenmessen und auf die Strahlenbrechung, von Hrn. Dr. \textsc{Benzenberg}. Das Urtheil über die zwar sinnreiche, aber bis jetzt doch noch zu problematische, Hypothese \textsc{Dalton}'s gehört vor das Forum des Physikers, und nicht vor das des Astronomen. Ob das, was Hr. \textsc{Benzenberg} auf Hrn. \textsc{Tralles} Einwürfe gegen jene Theorie erwidert, zulässig sei oder nicht, müsste doch, däucht uns, die Chemie leicht entscheiden können. Vor der Hand aber wäre es zu voreilig, jene Hypothese als ausgemachte Wahrheit bei den barometrischen Höhenmessungen zum Grunde zu legen. 24) Zufällige Gedanken über die Oberfläche des Mondes, von Hrn. Lieutenant \textsc{von Boguslawski}. Obgleich wir auf der Oberfläche des Mondes kein Wasser wahrnehmen, so könnte es doch, nach Hrn. \textsc{von Boguslawski}'s Meinung, dort vorhanden sein, nur auf der Tagseite von der Sonnenwärme in unsichtbare luftförmige Dämpfe verwandelt, die nach dem Untergang der Sonne nach und nach, wie sie die Wärme verlieren, in Nebel, Thau und Eis übergehen. Gegen diese Hypothese scheint uns das augenblickliche Verschwinden der Fixsterne, wenn sie vom dunklen Mondsrande bedeckt werden, ein wichtiger Einwurf zu sein. 25) Beobachtung der Pallas und Juno, Berechnung ihrer Gegenscheine, die Elemente der Bahn des letztern u.\,s.\,w. von Hrn. Prof. \textsc{Gauss}. Die beiden folgenden Artikel geben die Ephemeriden der Pallas und Juno für 1813 von den Herren \textsc{Nicolai} und \textsc{Wachter}. 28) Noch etwas über den wandelbaren Doppelstern 61 Schwan, Beobachtung des grossen Cometen von 1811, Sternbedeckungen und astronomische Nachrichten, von Hrn. Prof. \textsc{Bessel}. 29) Astronomische Beobachtungen und Bemerkungen, von Hrn. Dr. \textsc{Koch} in Danzig. 30) Ueber den neuen Cometen vom Jahre 1812. Der letzte Artikel gibt unter der Aufschrift: Vermischte astronomische Nachrichten, noch Beobachtungen des grossen Cometen von 1811 zu Petersburg von Hrn. \textsc{Schubert}; die Preise der astronomischen Instrumente, welche Hr. \textsc{Baumann} in Stuttgart verfertigt, und verschiedene andere astronomische Notizen. --- Wir schliessen diese Anzeige des an Reichhaltigkeit und Interesse sich immer gleich bleibenden Jahrbuches mit dem Wunsche, dass der würdige Herausgeber auch künftig bei seiner mühsamen Unternehmung kräftig unterstützt werden und uns noch lange alljährlich mit der Fortsetzung dieses Jahrbuches erfreuen möge.
